%%%%

%%%   Самарский государственный технический университет

%%%   Кафедра прикладной математики и информатики

%%%

%%%   Преамбула для статьи в журнал "Вестник Самарского государственного

%%%   технического университета. Серия: «Физико-математические науки»"

%%%   Ориентирована под MikTEx 2.4 for Win32

%%%

%%%   Хотя сам журнал собирается под GNU/Linux на дистрибутиве TeXLive



\documentclass[11pt,a4paper,twoside]{article}



\usepackage[cp1251]{inputenc}

\usepackage[T2A]{fontenc}

\usepackage[english,russian]{babel}

\usepackage{samgtubib}

\usepackage[dvips]{graphicx}

\usepackage[multiple]{footmisc}



\usepackage{samgtu}

\renewcommand{\VestnikYear}{2009} % Год издания Вестника

\renewcommand{\VestnikNumber}{2\,(19)} % Номер Вестника

\renewcommand{\VestnikMonth}{Ноябрь} % Месяц выхода Вестника

\renewcommand{\VestnikData}{01.11.09} % Дата подписания Вестника в печать

\renewcommand{\VestnikUPL}{26,64} % Усл. печ. л.

\renewcommand{\VestnikUIL}{25,73} % Уч.-изд. л.

\renewcommand{\VestnikRegNumber}{479/09} % Регистрационный номер

\renewcommand{\VestnikOrder}{994} % Номер заказа







%%

%

%\usepackage{pst-barcode}

%\usepackage{auto-pst-pdf}

%

%%\samgtupubl % äëÿ ãåíåðàöèè pdf, oòäàâàåìoão â òèïoãðàôèþ ÑàìÃÒÓ

%\pdfcrop % oáðåçêà pdf äëÿ ðàçìåùåíèÿ íà math-net.ru

%\draft % Íà êoððåêòóðå

%\filllastpage % Ïoñëåäíÿÿ ñòðàíèöà íå çàïoëíåíà íå ïoëíoñò

%

%\begin{document}

%



\normalsize





\vsgtu{1440}



\FirstPage{407}

\LastPage{414}





%







\renewcommand{\Vestnik}{% Êoìàíäà íà ïåðâûé âåðõíèé êoëoíòèòóë ñòàòüè

{\center

\parbox[c][7mm][s]{\textwidth}{\renewcommand{\bfdefault}{bx}\scriptsize \bf  Vestn.\ Samar.\ Gos.\ Techn.\ Un-ta. Ser. Fiz.-mat.\ nauki  %\  \VestnikYear. \ Issue~\VestnikNumber. \ Pp.\,\firstpage--\lastpage

\vfill

[J.~Samara State Tech. Univ.,\, Ser. Phys. \& Math. Sci.],   \VestnikYear,  vol.~\VestnikVolume,  no.~\VestnikNumber,   pp.\,\firstpage--\lastpage

\vfill 

\rule{\textwidth}{0.7pt}

\vfill 

\issnline \hfill \midoiline 

\renewcommand{\bfdefault}{b}

}}}









\chead[\fancyplain{}{\small\sl Oleg Igorevich Marichev (On the Occasion of His 70th Birthday)}]{\fancyplain{}{\small\sl Oleg Igorevich Marichev (On the Occasion of His 70th Birthday)}}

  \rhead[]{\fancyplain{\Vestnik}{}}

  \lhead[\fancyplain{\Vestnik}{}]{}

%  \rhead[\fancyplain{\Vestnik}{}]{\fancyplain{\thepage}{\thepage}}

%  \lhead[\fancyplain{\thepage}{\thepage}]{\fancyplain{\Vestnik}{}}

  \cfoot[]{}

  \lfoot[\thepage]{}

  \rfoot[]{\thepage}

  \plainheadrulewidth=0.7pt

  \headrulewidth=0.4pt



  \pagestyle{fancyplain}



  \thispagestyle{plain}



\addcontentsline{etoc}{section}{\it Oleg Igorevich Marichev (On the Occasion of His 70{\small th} Birthday)\smallskip}





\addcontentsline{toc}{section}{\it Îëåã Èãîðåâè÷ Ìàðè÷åâ (ê 70-ëåòèþ ñî äíÿ ðîæäåíèÿ)\smallskip}













\renewcommand{\baselinestretch}{0.95}

\normalfont





$\;$



\bigskip



\hfill \begin{pspicture}(1in,1in)

\psbarcode{http://dx.doi.org/10.14498/vsgtu1440}{}{qrcode}

\end{pspicture}

\vspace{-1in}





\entitle{Oleg Igorevich Marichev (On the Occasion of His 70{\small th} Birthday)}

\engtitlehackk{OLEG IGOREVICH MARICHEV \\ (On the Occasion of His 70{\small th} Birthday)}



\vskip8mm

\engTitleInTitlepage



\footnotetext{

\footnotesize\English

\vspace{-1em}\\

%\noindent\issnline; \midoiline\\ 

\textcopyright\ \VestnikYear\ Samara State Technical University. \vspace{.3em}\\ 

The electronic version of this paper contains hyperlinks to relevant websites. \vspace{.3em}\\

%\EnCitation 

\textbf{Please cite this article in press as:\\} {\engAuthorInContents} {\engTitleInContents}, \textit{Vestn. Samar. Gos. Tekhn. Univ., Ser. Fiz.-Mat. Nauki}\/ [J.~Samara State Tech. Univ., Ser. Phys. \& Math. Sci.],  \VestnikYear, vol.~\VestnikVolume, no.~\VestnikNumber,  pp.~\firstpage--\lastpage. \doihacke doi:~\href{http://mi.mathnet.ru/eng/\doisuffix}{\tt 10.14498/\doisuffix}.

\vspace{.3em} 

}



\vskip12mm



\begin{wrapfigure}{l}{0.46\textwidth} 

\centering

\includegraphics[width=0.4\textwidth]{oleg-marichev}



\end{wrapfigure} 



\English



\small



\smallskip



\noindent

\textbf{Early Years: 1945–1968}



\noindent

Oleg Igorevich Marichev was born on September 7, 1945 in Velikie Luki, Russia—a small town located near Pskov. Oleg's young parents Igor (23) and Lisa (21) were both students. When Oleg was five years old, Igor Marichev left the family to remarry, after which he received a Ph.D. in engineering, became a professor, and published several books. Lisa Marichev finished her studies and worked as an economist. Lisa Marichev died in 1983 and Igor passed away in 1998.



Oleg lived together with his mother and her sister Tamara (who later became a chemical researcher in the Belarus Academy of Sciences) in the extended family of his grandparents. Oleg's grandparents came from farming stock, but his grandfather Vasilii achieved the position of co-chief of police, while his grandmother Sofia kept house for the family. Vasilii and Sofia both died in 1982–1983, and Tamara passed away in 2001.



In 1949, the five members of the extended Marichev family relocated to Minsk, Belarus—a city that had been 

\href{https://us.search.yahoo.com/search?p=Minsk+photo+1944}{completely destroyed}

in the Second World War. During the immediate postwar period, herds of cows freely roamed the main streets of Minsk, and the Marichev family lived at one end of a long house that was still in the process of being completed by German prisoners of war. In addition to having difficult living conditions, Minsk was also a dangerous place for children, as a number of Oleg's friends were tragically killed by explosions after picking up and playing with live bombs that remained scattered about the city.



Between 1952–1963, Oleg attended school. His mathematical interest was first stimulated when his eighth grade mathematics teacher Alexandra Bagreeva explained the method of mathematical induction. Pursuing his new found interest in mathematics, Oleg became an avid participant in mathematical olympiads and competitions in Minsk. During this period, he attended a special mathematics school at 

\href{http://www.bsu.by/en/main.aspx}{Belorussian State University}, 

where he won many mathematical competitions. However, the contests also led in one very different direction, as Oleg met his future wife Anna in the course of competing in them. In the years 1960–1963, Oleg also pursued athletic interests as a member of the Belorussian youth triple and long jump teams, doing well enough to win second prize in a number of competitions.



In 1963, Oleg completed high school, from which he graduated with a prestigious academic gold medal—an honor bestowed upon only one other member of Oleg's graduating class of 100 students. Oleg then began undergraduate studies in the mathematics department at the Belorussian State University, graduating in 1968. Despite a slightly disappointing performance in ``scientific communism'' class (the one course in which he did not receive an ``excellent'' grade), upon graduation, Oleg was recommended to continue on with postgraduate studies in mathematics. 





\smallskip









\footnotetext{

\\

\noindent

\textbf{ADDITION (Information from Wiki)}





\noindent

Wiki gives the following information about \href{https://en.wikipedia.org/wiki/Oleg_Marichev}{Oleg}, which allows to see his scientific genealogy as following: Oleg defended Ph.D. in 1973 and is the student of \href{http://www.genealogy.ams.org/id.php?id=63083}{Fedor Dimitrievich Gahov}, who defended Ph.D. in 1937 and is the student of \href{http://www.genealogy.ams.org/id.php?id=35197}{B.~M.~Gagaev}, who defended Ph.D. in 1936 and is the student of \href{http://www.genealogy.ams.org/id.php?id=98816}{N.~N.~Parfent'ev}, who defended Ph.D. in 1911 and is the student of \href{http://www.genealogy.ams.org/id.php?id=105436}{A.~V.~Vasil'ev}, who defended Ph.D. in 1884 and is the student of \href{http://www.genealogy.ams.org/id.php?id=148745}{P.~L.~Chebyshev}, who defended Ph.D. in 1849 and is the student of \href{http://www.genealogy.ams.org/id.php?id=12542}{N.~D.~Brashman}, who defended Ph.D. in 1834 and is the student of \href{http://www.genealogy.ams.org/id.php?id=12541}{N.~I.~Lobachevsky},  who is the student of

\href{http://www.genealogy.ams.org/id.php?id=146776}{J.~M.~Ch.~Bartels}.

Then genealogy tree of mathematicians matches \href{http://www.genealogy.ams.org/id.php?id=18231}{C.~Gau\ss (1799)}, \href{http://www.genealogy.ams.org/id.php?id=35953}{A.~M\"oebius (1815)}, \dots, \href{http://www.genealogy.ams.org/id.php?id=60985}{G.~Leibniz (1666)} and others.



\indent 

Short information about Oleg can be found at \href{http://functions.wolfram.com/About/developers.html}{the Developers of the Mathematical Functions Website}. 

The picture of his books and Functions poster can be found at the \href{http://www.stephenwolfram.com/publications/history-future-special-functions/intro.html}{Stephen Wolfram' Official Website}.

The article \href{http://web.archive.org/web/20120212191344/http://www.wolfram.com/news/events/techconf2005/marichev.html}{Special Session: In Honor of Oleg Marichev on the Occasion of His 60th Birthday}

includes the \href{http://web.archive.org/web/20120922034933/http://www.wolfram.com/news/events/techconf2005/MarichevPublications.nb}{list of his publications}. 

Additionally to that list he has published monograph about partial differential equations

\href{https://www.zbmath.org/?q=an:1147.35003}{Marichev,~O.~I.; Kilbas,~A.~A.; Repin,~O.~A. (2008) \textit{Boundary value problems for partial differential equations with discontinuous coefficients} (Russian)}  and several

electronic blogs (see \url{http://blog.wolfram.com/author/oleg-marichev/}).

}







\noindent

\textbf{Academic Years: 1968–1990}



\noindent

In 1968, Oleg began his research mathematics career under the guidance of academician and Academy of Sciences of Belarus Professor \href{http://www.genealogy.ams.org/id.php?id=35197}{Fedor Dimitrievich Gakhov}. Gakhov was famous for having solved the classical Riemann boundary value problem for analytic functions in closed form in 1937, for his monograph \href{http://www.worldcat.org/title/boundary-value-problems/oclc/1722840&referer=brief_results}{\textit{Boundary Value Problems}} that was translated into English in 1966, and for his supervision of nearly 60 Ph.D. dissertations. Gakhov recognized Oleg's talent and ability to work independently, and suggested a thesis topic that strayed far from his own research while still building on some of his results. Following Gakhov's advice, Oleg started research on mixed-type partial differential equations (i.e., elliptic-hyperbolic equations). While investigating Tricomi-type boundary value problems for these equations, Oleg discovered that such problems could be converted into singular integral equations that could then be solved using Gakhov's algorithms for Riemann boundary problems. Oleg also constructed solutions for several classes of those equations, as well as for the related Euler–Poisson–Darboux generalized equations.



While engaged in his thesis research in 1970, Oleg first encountered hypergeometric, Bessel $J$-, Legendre $P$-, Appell $V$-, and Meijer $G$-function, thus beginning a lifelong fascination with special functions. In particular, Oleg's interest was piqued by a copy of Lucy Joan Slater's book \href{http://www.amazon.co.uk/Generalized-Hypergeometric-Functions-Lucy-Slater/dp/052109061X}{\textit{Generalized Hypergeometric Functions}} that he had borrowed from Lenin's Library in Moscow (and a copy of which he subsequently received from Slater as a gift). Upon perusing its well-written pages laying out a map of the hypergeometric world, Oleg came to recognize the beautiful applications and problems that could be solved using these functions. Applying integral equations and special functions to the solutions of partial differential equations became his thesis, and Oleg defended his \href{http://www.genealogy.ams.org/id.php?id=63083}{(first) Ph.D. in 1973}. (In the former Soviet Union, it was possible to obtain two doctorate degrees, the first being the usual Western Ph.D., and the second being an ``advanced Ph.D.'', or Doctor Sciencarum, a special degree awarded for outstanding and prolific work in a given field. Oleg is the holder of both, having been awarded the latter from Friedrich--Schiller--Universit\"at in 1990.)



Oleg had begun teaching at the Belorussian State University Mathematics Department starting in 1968, and continued doing so until 1990. One of his first students (in 1969–1970) was \href{http://www.mathnet.ru/eng/person48302}{Sergei Rogozin}, who is now a professor in that department. Ivan Sandrigailo, another of Oleg's students, won first place in Soviet student mathematical competitions, but tragically died in 1973.



The early years of teaching and research also brought sorrow to Oleg's own family. He had married his first wife Lubov in 1969, and the two conceived a son who was to be born in 1972. After carrying the child for nine months, tragedy struck, as emergency medical procedures proved insufficient to save the lives of either Lubov or the unborn son she carried. The deaths of both wife and child were a hard blow for Oleg, as well as for everyone who knew Lubov. After several months of mourning and poor physical health, Oleg married his former mathematical competitor Anna, and the two joined their lives in a lifelong partnership that has thus far spanned more than 42 years.



While teaching, Oleg began authoring a series of authoritative (and massive) monographs about special functions. The first of these, \textit{Method vychisleniya integralov ot spetsialnykh funktsii (teoriya i tablitsy formul)}, was published in Russian in 1978 and was subsequently translated into English as 

\href{http://www.worldcat.org/title/handbook-of-integral-transforms-of-higher-transcendental-functions-theory-and-algorithmic-tables/oclc/8866416}{\textit{Handbook of Integral Transforms of Higher Transcendental Functions} (Ellis Horwood Ltd., 1983)}. This book characterized the completeness and rigor in Oleg's publications, developing and extending ideas first considered in the 1940s by \href{https://en.wikipedia.org/wiki/Cornelis_Simon_Meijer}{C.~S.~Meijer} in a way that was characterized by none other than Gakhov himself as, ``[without] parallel in the Russian or foreign literature''. \textit{Handbook of Integral Transforms} included the largest table of the 

\href{http://mathworld.wolfram.com/MeijerG-Function.html}{Meijer $G$-function} ever compiled, a feat of considerable interest because it allowed most integrals tabulated in handbooks such as Gradshteyn and Ryzhik's famous 

\href{http://www.mathtable.com/gr/}{\textit{Tables of Integrals, Series, and Products}}

to be immediately evaluated as particular cases of the \href{http://mathworld.wolfram.com/MeijerG-Function.html}{Meijer $G$-function}.



(Ironically, \textit{Handbook of Integral Transforms} was actually Oleg's second monograph penned in 1975–1977. His first book, \textit{Boundary Value Problems for Partial Differential Equations with Discontinuous Coefficients}, could not be published in 1977 on the grounds that Oleg was too young to publish two books simultaneously and therefore must postpone publication of one of them until the following year. 

Regrettably, this first book remains unpublished and, 

only 31 years later, it was published with co-authors \href{https://www.zbmath.org/?q=an:1147.35003}{A.~A.~Kilbas  and O.~A.~Repin}.)



After publishing \textit{Handbook of Integral Transforms}, Oleg got in touch with Moscow professor 

\href{https://en.wikipedia.org/wiki/Anatolii_Platonovich_Prudnikov}{Anatoli Prudnikov}

at the suggestion of Gakhov. Prudnikov proposed a collaboration with himself and 

\href{https://www.researchgate.net/profile/Yury_Brychkov2/topics}{Yury Brychkov}

for the purposes of compiling a set of integral tables that would be even more comprehensive than Gradshteyn and Ryzhik's, a proposal that Oleg readily accepted. Oleg thus began the huge task of evaluating and re-evaluating many thousands of complicated integrals and other formulas, all without the benefit of even rudimentary computers. Brychkov worked mostly on series and indefinite integrals, and Prudnikov kept the project moving, including the nontrivial task of assuring publication despite various administrative and bureaucratic difficulties. While engaged in this gargantuan task, Oleg calculated integrals around the clock and all about town—during meetings, on the metro or trolley, and before falling asleep. As result of these determined efforts, \href{https://www.zbmath.org/?q=an:1103.00302|an:0606.33001|an:0511.00044}{volumes 1–3} of 

\href{https://www.zbmath.org/?q=ai:%22marichev.oleg-i%22%20ti:%22Integrals%20and%20Series%22}{\textit{Integrals and Series}} 

were published in Russian and 

\href{https://www.zbmath.org/?q=an:0781.44002|an:0786.44003}{volumes 4–5

in English} during the eleven-year period 1981–1992. (The first two volumes have now also been published in Japanese, and Russian reprints of \href{https://www.zbmath.org/?q=an:1103.00300|an:1103.00301|an:1103.00302}{volumes 1–3 appeared in 2003}, together with a small student version entitled 

\href{http://www.worldcat.org/title/tables-of-indefinite-integrals/oclc/19122709&referer=brief_results}{\textit{Tables of Indefinite Integrals}}. Picture of all Oleg's book see \href{http://www.stephenwolfram.com/publications/history-future-special-functions/intro.html}{here}.)



In 1975, Oleg and colleague 

\href{http://academic.research.microsoft.com/Author/12834930/anatoly-a-kilbas}{Anatoli Kilbas}

traveled to Odessa to attend a conference. While there, they decided to write a book on fractional calculus and sought to enlist the collaboration of fractional differentiation expert 

\href{http://www.genealogy.ams.org/id.php?id=106361}{Stephan Samko}

(who also happened to be a former Ph.D. student of Gakhov). Samko agreed to join the project, and the three coauthors enthusiastically set out on this task that they completed several years later in the form of the monograph

\href{http://www.amazon.com/Fractional-Integrals-Derivatives-Theory-Applications/dp/2881248640}{\textit{Fractional Integrals and Derivatives (Theory and Applications)}}, published in Minsk in 1987 and in an expanded English version by Gordon and Breach in 1993. 



The years 1981–1990 were a period of intensive scientific work by Oleg and a group of six of his Ph.D. students (Nguyen Ti Than, 

\href{http://academic.research.microsoft.com/Author/12980097/vu-kim-tuan}{Vu Kim Tuan}, 

\href{http://www.cs.cmu.edu/~adamchik/}{Victor Adamchik}, 

\href{http://academic.research.microsoft.com/Author/54570726/semen-b-yakubovich}{Semen Yakubovich}, 

Galina Grinkevich, and 

\href{http://www.goodreads.com/author/show/573284.Thanh_Hai_Nguyen}{Nguyen Thanh Hai}). Oleg's students defended doctoral dissertations in 1986–1988, and Vu Kim Tuan even defended a second Ph.D. in 1987 (after having published nearly 20 articles in the years 1981–1985). In 1990, Oleg himself received a second Ph.D. from Friedrich–Schiller–Universit\"at in Jena, Germany—the institution of his co-author 

\href{http://www.mathnet.ru/eng/person38690}{H.-J. Glaeske}. 

Oleg's second Ph.D. was especially unusual because it represented the first time a second Ph.D. in mathematics had ever been defended outside the Soviet Union.



In the midst of this flurry of academic activity, Oleg and Anna's son Sergei was born in 1983. In 1986, the Chernobyl nuclear reactor accident and explosion had a huge effect on many people, including Oleg's family. Oleg himself was hospitalized twice for a severe flu that his radiation-weakened immune system was not capable of fighting off by itself.



As early as 1978, Oleg had dreamed of automating the algorithm for evaluation of integrals that he had developed in his first book. Even though computer technology in the Soviet Union lagged behind that in the West, Oleg collaborated with

\href{http://www.mathnet.ru/eng/person114502}{Ernst Krupnikov} 

in Novosibirsk in 1980 on the MIR computer. Oleg and Krupnikov where able to implement integration using the convolution theorem with \href{http://mathworld.wolfram.com/MeijerG-Function.html}{Meijer $G$-function} and, in 1987, Oleg suggested further work along these lines to his former Ph.D. student Victor Adamchik. Adamchik worked in a new departmental computer laboratory established by Oleg and A.~Kilbas, and implemented a prototype integration system using the language REDUCE. While reporting on this work at a computer conference in Germany in 1990, Oleg heard that

\href{http://www.stephenwolfram.com/}{Stephen Wolfram}, 

founder of a new technology company known as \href{http://www.wolfram.com/}{Wolfram Research}, was becoming interested in effective algorithms for evaluation integrals. Very shortly thereafter in May of 1990, Oleg received a telephone invitation from \href{http://www.wolfram.com/}{Wolfram Research} to come to the United States with Adamchik and demonstrate their program. Following discussions with then Wolfram employee 

\href{http://blog.wolfram.com/author/roman-maeder/}{Roman Maeder}

at the August 1990 SIAM conference in Japan, Oleg agreed to come to the company's headquarters in Champaign at the end of 1990.



Since 1987, Oleg and Anna had started to think about looking for jobs at a university in the West. There were many reasons for thinking this, not only because of Chernobyl-induced health issues and worries, but also because of the overall unraveling of the institutional fabric of the Soviet Union. From a technological standpoint, easy access to advanced computer technology in the West meant implementing a grand integration algorithm would also be much easier there. There was also a very different but equally important reason: official corruption and brutality within the Soviet government. Shortly after Oleg and his colleagues Prudnikov, Brychkov, 

\href{www.goodreads.com/author/show/6143782.Mikhail_V_Fedoryuk}{Mikhail Fedoryuk}, and 

\href{http://www.worldcat.org/identities/lccn-n88-9291/}{\`Eduard Riekstyn'sh} were presented in 1989–1990 with the highest Soviet award for their series of books about integrals, professor and coauthor Fedoryuk was beaten to death by officials of State Security following an altercation with a ticket collector at the Kiev railway station as he returned home to Moscow. The fact that a mathematics professor could be killed for no reason and that the somebody responsible was never even brought to trial for this crime was the final straw that led Oleg and his family to decide to immigrate to the United States in 1990–1991.





\pagebreak



\noindent

\textbf{Wolfram Years: 1990–}



\noindent

In November 1990, Oleg and Adamchik arrived in Champaign. After Oleg and Adamchik demonstrated their integration program, Stephen Wolfram suggested to extend their visit by three months in order to experiment with implementing their algorithms in \textit{Mathematica}. They continued to work on this project, and they both received temporary visas over the next few months. Oleg's wife Anna and son Sergei (8 years old at that time) joined Oleg in Champaign in March 1991, and Oleg received permanent resident status in 1992. Anna started to work as a volunteer and also began the process of learning English. She worked at \href{http://www.wolfram.com/}{Wolfram Research} for six years, moved on to the integrated customer management firm AMDOCS, and now works at the University of Illinois.



The first part of implementing integration in \textit{Mathematica} concerned the indefinite integration of elliptic integrals, and was completed by Oleg, Adamchik, and 

\href{http://research.microsoft.com/en-us/um/people/alexeib/}{Alexei Bocharov}, with some help from Emily Martin.  This part was realized on Mathwematica and continued C-code for Risch algorithm, that was written by

\href{http://kellyroach.com/index.htm}{Kelly Roach}.

 The next part consisted of indefinite integrals that could be expressed in terms of special functions. After completion of this work, \textit{Mathematica}’s \href{http://reference.wolfram.com/language/ref/Integrate.html}{\texttt{Integrate}} function was now capable of evaluating essentially all indefinite integrals from Oleg's tables. Simultaneously, Oleg, Adamchik, and Bocharov implemented definite integration, and the three continued to improve and extend \href{http://reference.wolfram.com/language/ref/Integrate.html}{\texttt{Integrate}}. After Bocharov and Adamchik departed from \href{http://www.wolfram.com/}{Wolfram Research}, this important and difficult work has been continued by \href{https://www.linkedin.com/in/daniellichtblau}{Daniel Lichtblau} and later 

\href{https://www.linkedin.com/in/oleksandr-pavlyk-39a4034}{Oleksandr Pavlyk}, 

\href{http://community.wolfram.com/web/dkapadia}{Devendra Kapadia} with assistance from Oleg.



In the period 1992–1997, Oleg actively worked not only on symbolic integration but, together with 

\href{http://www.stephenwolfram.com/publications/jerry-keiper/}{Jerry Keiper} 

(who died in a tragic accident in early 1995), also on numerical evaluation the \href{http://mathworld.wolfram.com/MeijerG-Function.html}{Meijer $G$-function}—the most complicated of \textit{Mathematica}'s special functions. Together with 

\href{http://members.wolfram.com/adams/}{Adam Strzebonski}, 

he also introduced new \href{http://reference.wolfram.com/language/ref/FunctionExpand.html}{\texttt{FunctionExpand}} functionality and, together with Lichtblau, improved some basic formulas for \href{http://reference.wolfram.com/language/ref/Solve.html}{\texttt{Solve}}.



At 1997, Oleg began a productive collaboration with colleague and special functions expert 

\href{http://community.wolfram.com/web/mtrott}{Michael Trott}. 

Together, Oleg and Trott began the ambitious project of describing and classifying all \textit{Mathematica}’s mathematical functions by means of formulas, graphics, and interrelations. This project began on paper, with the original intent to create a set of (eventually) five mathematical posters cataloging the results. However, as the mass of accumulated material quickly outpaced the space available (even on five large printed posters), the cost of actually producing and distributing such complicated documents also spiraled upward, and the project moved to the internet, where it eventually became \href{http://functions.wolfram.com/}{The Wolfram Functions Site}. This website is now the world's most comprehensive collection of mathematical formulas, and it currently includes approximately 307,000 formulas and graphics, as well as a unique search functionality for formulas based on their mathematical structure. This work was given an additional boost through a \href{https://nsdl.oercommons.org/}{National Science Foundation Digital Library} grant obtained by Trott

%and colleague Eric Weisstein at \href{http://www.wolfram.com/}{Wolfram Research} together 

and colleague \href{http://www.ericweisstein.com/}{Eric Weisstein} (the author of famous \href{http://mathworld.wolfram.com/}{mathworld}) at \href{http://www.wolfram.com/}{Wolfram Research} together  

 with colleagues from the mathematics and engineering libraries at the University of Illinois.







Oleg's work on special functions continues not only through collaboration with Trott on \href{http://functions.wolfram.com/}{The Wolfram Functions Site}, but also with new colleague and special functions expert Olexander Pavlyk, who is now making it possible to further extend the functionality in \textit{Mathematica} functions such as \href{http://reference.wolfram.com/language/ref/MeijerG.html}{\texttt{MeijerG}}, \href{http://reference.wolfram.com/language/ref/FunctionExpand.html}{\texttt{FunctionExpand}}, and so on. Oleg also continues to work with Lichtblau on the extension of \href{http://reference.wolfram.com/language/ref/Integrate.html}{\texttt{Integrate}} and \href{http://reference.wolfram.com/language/ref/Series.html}{\texttt{Series}} and with Strzebonski on new capabilities for \href{http://reference.wolfram.com/language/ref/FunctionExpand.html}{\texttt{FunctionExpand}} and 

\href{http://reference.wolfram.com/language/ref/PowerExpand.html}{\texttt{PowerExpand}}.  \href{http://functions.wolfram.com}{The Wolfram Functions Site} acts as a focal point for all these labors, allowing Oleg to realize his dream of building a ``super handbook'' of formulas for special and elementary functions, as complete as the five volumes of \textit{Integrals arid Series}—and implemented entirely within \textit{Mathematica}. Oleg looks forward to pursuing this goal over the next 20 years together with friends and colleagues at \href{http://www.wolfram.com/}{Wolfram Research}.











During the last ten years Oleg's goal has begun its realization in several directions. Users can now find 307,000+ formulas at the \href{http://functions.wolfram.com}{Wolfram Functions Site} (see \href{http://blog.wolfram.com/2008/05/06/two-hundred-thousand-new-formulas-on-the-web}{Wolfram Blog}), 35,000+ formulas for 500+ probability distributions as shown in the \href{http://blog.wolfram.com/2013/02/01/the-ultimate-univariate-probability-distribution-explorer}{Wolfram blog}, about 10,000 formulas for continued fractions through the \href{http://www.wolframalpha.com}{WolframAlpha website}, as mentioned in the \href{http://blog.wolframalpha.com/2013/05/01/after-100-years-ramanujan-gap-filled/}{WolframAlpha  Blog}, and implementation of the majority of formulas from the \href{http://functions.wolfram.com}{Wolfram Functions Site} into the 

%\textit{Mathematica} system (Version 10.3) through the use of the Wolfram language 

\textit{Mathematica} system (Version 10.3, implemented by \href{https://www.linkedin.com/in/paco-jain-20472297}{Paco Jain}) through the use of the Wolfram language    

symbols \href{https://reference.wolfram.com/language/ref/MathematicalFunctionData.html}{\texttt{MathematicalFunctionData}} and

\href{https://reference.wolfram.com/language/ref/EntityValue.html}{\texttt{EntityValue}}.





Other results are not visible yet but they are approaching a good stage of implementation. They are the following:

\itemsep=0mm  

\parsep=0mm 

 \begin{itemize}

 \small

\itemsep=0mm  

\parsep=0mm 

\item[1)] improving \textit{Mathematica}'s \href{http://reference.wolfram.com/language/ref/Integrate.html}{\texttt{Integrate}} on the through the use of the \href{http://reference.wolfram.com/language/ref/MeijerG.html}{\texttt{MeijerG}} function, with explanations of how to evaluate integrals 

via \href{http://www.wolframalpha.com/pro/step-by-step-math-solver.html?src=widget}{Step-by-Step technology} 

%(see also \href{http://integrals.wolfram.com/index.jsp}{Wolfram Mathematica Online Integrator}), 

(see also \href{http://integrals.wolfram.com/index.jsp}{Wolfram Mathematica Online Integrator}), project is realising with professor \href{http://www.cs.cmu.edu/~adamchik/}{Victor Adamchik} and \href{http://blog.wolframalpha.com/author/ghurst/}{Greg Hurst},



\item[2)]  incorporating near 500 mathematical constants and hundreds of new functions with their definitions and formulas into \href{http://reference.wolfram.com/language/ref/MathematicalFunctionData.html}{\tt MathematicalFunctionData}, 



\item[3)]  separations of Laplace, Helmholtz and other partial differential equations on the 50+ different 

%coordinate systems,

coordinate systems (with \href{https://www.cs.hmc.edu/~is/}{Itai Seggev}),



\item[4)]  building correct theory of $q$-\href{http://reference.wolfram.com/language/ref/MeijerG.html}{\texttt{MeijerG}} with its hundred cases (with \href{https://www.researchgate.net/profile/Yury_Brychkov2/topics}{Yury Brychkov}), 



\item[5)]   Implementation effective program for evaluation continued fractions and symbolical derivatives (with \href{http://www.math.buffalo.edu/mad/PEEPS2/pooh_charles.html}{Charles Pooh}),





\item[6)]  developing of different types of converters for mathematical functions and their visualization on Internet using modern \textit{Mathematica} technology (prototype of this idea was realized near 18 years ago in \href{http://www.cybertester.com/data/mtposter.pdf}{The Poster  

\textit{The Mathematical Functions of Mathematica}},  see also  \href{http://www.mathematica25.com/}{The Mathematica Story: A~Scrapbook}).



\end{itemize}









All formulas, which Oleg finds in literature or derives himself, go through several stages: transformation to \textit{Mathematica}'s format, verification, possible extention to wider domain with minimum restrictions and conversion to \href{https://reference.wolfram.com/language/ref/TraditionalForm.html}{\tt Traditional\-Form}. 

Only after these steps the formula can be classified and placed into the corresponding subsection of corresponding file or other document.





In all of the above activity Oleg tries to realize a general axiomatic approach: in each area (functions, integrals, probability distributions, continued fractions, coordinate systems, differential equations, etc.) he finds and works with maximally general or abstract functions and after operation with corresponding properties he drills down to concrete cases. This approach was realized in realms of function identities and integration where he uses the \href{http://reference.wolfram.com/language/ref/MeijerG.html}{\texttt{MeijerG}}  (which he himself implemented) and the more general \href{http://mathworld.wolfram.com/FoxH-Function.html}{\tt FoxH} (yet to be implemented, hopefully even with two variables) functions, in probability distributions (where he derived general formulas for the Pearson distribution with six parameters and described properties of abstract distributions), and in differential equations connected with ellipsoidal coordinate systems (where Oleg found integral representation through Lauricella functions for general Heun type differential equation with four regular singular points, which generalizes famous Riemann differential equation with three regular singular points for Gauss hypergeometric function ${}_2F_1$).



But the largest revolutional change was made in the all mathematical formulas, involving multivalued functions with branch cuts (including approximately half of the elementary functions). During centuries people used formulas like $\sqrt{1-z^2}= \sqrt{1-z} \sqrt{1+z}$ and tried not to use $\sqrt{z^2-1}=\sqrt{z-1}\sqrt{z+1}$ (the first is correct for all complex values of $z$, but the second is not correct for $z<-1$, where the right part must have additional factor $-1$, implementation of which allows to make the following correct everywhere formula 

$$\sqrt{z^2-1}\equiv \sqrt{z-1} \sqrt{z+1} \exp\Bigl(i \pi \Bigl\lfloor \frac{1-\arg(z-1)-\arg(z+1)}{2 \pi}

\Bigr\rfloor\Bigr), $$ 

where $\lfloor \text{\textvisiblespace} \rfloor$ is the floor function.

During 20+ years in \href{http://www.wolfram.com/}{Wolfram Research Inc.}  Oleg undertook the huge job of converting existing, and deriving new formulas for all mathematical functions to be consistent with the general axiomatical restriction: ``The argument of all complex numbers $z$ satisfies the  inequality 

${-\pi < \arg(z)\leq \pi}$''. 

Appearance of computers (which can operate now only on one complex plane) demanded such precision, which led to the appearance of a new ``computer's world''  of mathematical formulas, presented in \href{http://functions.wolfram.com/}{Wolfram Functions Site} and supported by {\it Mathematica} and other systems. 

In particular, 

%very basic formulas for logarithm of product, and power expanding of products, were corrected (see

%\href{http://functions.wolfram.com/ElementaryFunctions/Log/16/04/01/0008/}{Natural logarithm: Transformations (formula 01.04.16.0018)}  and  \href{http://functions.wolfram.com/ElementaryFunctions/Power/16/04/01/01}{Power function: Transformations (subsection 16/04/01/01)}). 

very basic formulas for logarithm of product, and power expanding of products, were corrected (see \href{http://functions.wolfram.com/ElementaryFunctions/Log/16/04/01/}{Natural logarithm: Transformations} (formulas \href{http://functions.wolfram.com/ElementaryFunctions/Log/16/04/01/0008/}{01.04.16.0018} and \href{http://functions.wolfram.com/ElementaryFunctions/Log/16/04/01/0009/}{01.04.16.0076}) and \href{http://functions.wolfram.com/ElementaryFunctions/Power/16/04/01/01/}{Power function: Transformations (9 formulas in subsection 16/04/01/01)}).    

Of couse, it effected majority formulas for more complicated elementary and special functions 

(see thousands new inter-relation formulas, connected all 26 direct and inverse trigonometric and hyperbolic, exponential and logaritm functions, and, for example, 275 formulas, 

that connected \href{http://reference.wolfram.com/language/ref/ArcSin.html}{\tt ArcSin} and \href{http://reference.wolfram.com/language/ref/ArcTan.html}{\tt ArcTan} functions in \href{http://functions.wolfram.com/ElementaryFunctions/ArcSin/27/02/03}{Inverse sine: Representations through equivalent functions (subsection 27/02/03)}).

Many classical functions and formulas were able to be defined for a wider portion of the complex plane  (for example, Stirling formula for Euler's gamma function was extended for negative arguments in \href{http://functions.wolfram.com/GammaBetaErf/Gamma/06/02/0005/}{Gamma function: Series representations (formula 06.05.06.0042)}, 

polygamma function was defined for arbitrary complex index

in \href{http://functions.wolfram.com/GammaBetaErf/PolyGamma2/02/}{Polygamma function: Primary definition}, asymptotical 

formula for Bessel functions $J_{\nu}(z)$ was 

extended on negative values $z$ (for 

${\arg (z) = \pi \land z \to -\infty}$) (see 35 formulas in \href{http://functions.wolfram.com/Bessel-TypeFunctions/BesselJ/06/02/}{Bessel function of the first kind: Series representations (subsection 06/02)}) 

and the idea of fractional integro-differentiation was developed more deeply, which allowed many new formulas for 

%symbolical and fractional differentiation to be established  (see  \href{http://functions.wolfram.com/ElementaryFunctions/Power/20/03/01/}{Power function: Differentiation (subsection 20/03/01)}, \href{http://functions.wolfram.com/ElementaryFunctions/Log/20/}{Natural logarithm: Differentiation}, \href{http://functions.wolfram.com/Notations/3/}{Listing of the Mathematical Notations used in the Mathematical Functions Website: Operations}). For deriving 

symbolical and fractional differentiation to be established (see \href{http://functions.wolfram.com/ElementaryFunctions/Power/20/03/01/}{Power function: Differentiation (subsection 20/03/01)}, 25 formulas in \href{http://functions.wolfram.com/ElementaryFunctions/Log/20/}{Natural logarithm: Differentiation}, many formulas listing of the  \href{http://functions.wolfram.com/Notations/3/}{Mathematical Notations used in the Mathematical Functions Website: Operations}). For deriving 

and demonstration many calculus related formulas Oleg 

%used generic \href{http://reference.wolfram.com/language/ref/BellY.html}{\tt BellY} polynomials, which gives uniformity 

used generic \href{http://reference.wolfram.com/language/ref/BellY.html}{\tt BellY} polynomials (see \href{http://reference.wolfram.com/language/ref/BellY.html}{Documentation Center for Wolfram Language \& System}), which gives uniformity 

for the presentation of composition and inversion of series and functions and their derivatives.





Currently Oleg works with orthogonal classical and $q$-polynomials, and particularly towards extending the famous \href{http://dlmf.nist.gov/search/search?q=Askey+Scheme}{Askey scheme} of orthogonal polynomials. The main formulas for these and other mathematical functions will eventually be included into \textit{Mathematica} and supported by the Wolfram language. 





Below one can find 11 main Oleg's books in chronological order. 

Two of them (1983---Ellis Horwood and 1993) are translations

of Russian versions of books from 1978 and 1987, 

but with essential extensions. 

Some of books were republished on Russian, Japanese, German and English languages.



\begin{enumerate}

\itemsep=0mm  

\parsep=0mm  \small



\item \href{https://www.zbmath.org/?q=an:0473.33001}{Marichev, O. I. (1978) \textit{Metod vychisleniya integralov ot spetsial'nykh funktsij (Teo\-riya i tablitsy formul)} [A method of calculating integrals of special functions. (Theory and tables of formulas)] (Russian), Nauka i Tekhnika, Minsk, 310 pp.}





\item 

\href{https://www.zbmath.org/?q=an:0511.00044}{Prudnikov, A.~P.; Brychkov, Yu.~A.; Marichev, O.~I. (1981) \textit{Integraly i ryady. Ehlementarnye funktsii} [Integrals and series. Elementary functions] (Russian), Nauka, Moscow, 800 pp.}





\item

\href{https://www.zbmath.org/?q=an:0626.00033}{Prudnikov, A.~P.; Brychkov, Yu.~A.; Marichev, O.~I. (1983)

\textit{Integraly i ryady. Spetsial'nye funktsii} [Integrals and series. Special functions] (Russian), Nauka, Moscow, 752~pp.}



\item

\href{https://www.zbmath.org/?q=an:0494.33001}{Marichev, O.~I. (1983)

\textit{Handbook of integral transforms of higher transcendental functions: theory and algorithmic tables}, Ellis Horwood Ltd. Chichester, New York, 336~pp.}



\item  

\href{https://www.zbmath.org/?q=an:0606.33001}{Prudnikov, A.~P.; Brychkov, Yu.~A.; Marichev, O.~I. (1986)

\textit{Integraly i ryady. Dopolnitel'nye glavy} [Integrals and series. Complementary chapters]  (Russian), Nauka, Moscow, 800~pp.}





\item 

\href{https://www.zbmath.org/?q=an:0626.00032}{Brychkov,~Yu.~A.; Marichev,~O.~I.; Prudnikov,~A.~P. (1986)

\textit{Tablitsy neopredelennykh integralov. Spravochnik} [Tables of indefinite integrals. Handbook] (Russian), Nauka, Moscow, 192~pp.}





\item  

\href{http://www.worldcat.org/title/integraly-i-proizvodnye-drobnogo-poriadka-i-nekotorye-ikh-prilozheniia/oclc/18496926&referer=brief_results}{Samko, S. G.; Kilbas, A. A.; Marichev, O. I. (1987) \textit{Integraly i proizvodnye drobnogo poryadka i nekotorye ikh prilozheniya} [Integrals and Derivatives of Fractional Order and Some of Their Applications] (Russian), Nauka i Tekhnika, Minsk, 688 pp.}



\item  

\href{https://www.zbmath.org/?q=an:0786.44003}{Prudnikov, A. P.; Brychkov, Yu. A.; Marichev, O. I. (1992) \textit{Integrals and series}. Vol. 4. Direct Laplace transforms, Gordon and Breach Science Publ., New York, NY, 619 pp.}



\item  

\href{https://www.zbmath.org/?q=an:0781.44002}{Prudnikov, A. P.; Brychkov, Yu. A.;  Marichev, O. I. (1992)  \textit{Integrals and series}. Vol. 5. Inverse Laplace transforms, Gordon and Breach Science Publ., Philadelphia, PA, 595 pp.}



\item 

\href{https://www.zbmath.org/?q=an:0818.26003}{Samko, S. G.; Kilbas, A. A.;  Marichev, O. I. (1993) \textit{Fractional integrals and derivatives. Theory and Applications}, Gordon and Breach Science Publ., New York, NY, 976 pp.}



\item 

\href{https://www.zbmath.org/?q=an:1147.35003}{Marichev, O. I.; Kilbas, A. A.; Repin, O.~A. (2008) \textit{Kraevye zadachi dlya uravnenii v chastnykh proizvodnykh s razryvnymi koeffitsientami} [Boundary value problems for partial differential equations with discontinuous coefficients] (Russian), Samara State Economic Univ., Samara, 275~pp.}





\end{enumerate}





\bigskip



\textit{The editorial  board and  editorial  staff  together  with  the  scientific  community  of 

the  Samara State Technical University  congratulate  Oleg Igorevich Marichev  on  his  70{\small th}  anniversary and wish him good health, new creative achievements, and successes in science and his multifaceted activity.}







\renewcommand{\baselinestretch}{0.9}





\renewcommand{\Vestnik}{% Êîìàíäà íà ïåðâûé âåðõíèé êîëîíòèòóë ñòàòüè

%\mbox{\renewcommand{\bfdefault}{bx}\scriptsize \bf  Âåñòí.\ Ñàì.\ ãîñ.\ òåõí.\ óí-òà. Ñåð. Ôèç.-ìàò.\ íàóêè.  \  \VestnikYear. \ \No~\VestnikNumber. \ Ñ.\,\firstpage--\lastpage 

%\issnline

%\renewcommand{\bfdefault}{b}}

\parbox[c][7mm][s]{\textwidth}{\renewcommand{\bfdefault}{bx}\scriptsize  \bf  Âåñòí.\ Ñàì.\ ãîñ.\ òåõí.\ óí-òà. Ñåð. Ôèç.-ìàò.\ íàóêè.  \  \VestnikYear. \ T.~\VestnikVolume, \No~\VestnikNumber. \ Ñ.\,\firstpage--\lastpage \vfill

\issnline \hfill \doiline

\renewcommand{\bfdefault}{b}}\vspace{-9mm} 

}





\newpage

\Russian





%\end{document}

